% SEGMENT OLP-0683-S01
% Part: proof-theory
% Chapter: proof-search
% Section: search-algorithm

\documentclass[../../../include/open-logic-section]{subfiles}

\begin{document}

\olfileid{pt}{ps}{sal}

\olsection{\Log{G3c} க்கான நிறுவல் தேடல் படிமுறை}

\Log{G3c} க்கான ஒரு நிறுவல் தேடல் படிமுறையை
விளக்குகிறோம். இதுவே சாத்தியமான ஒரே படிமுறை அல்ல;
மிகவும் திறமையான படிமுறையும் அல்ல. ஆனால் இது சரியானது:
$\Gamma \Sequent \Delta$ என்ற தொடரணிக்கு
!!a{proof} இருந்தால், படிமுறை இறுதியில் நின்று
!!a{proof} ஐக் கண்டுபிடிக்கும். இது நிறுவலைத்
தவறவிடாததும் ஆகும்: !!a{proof} ஐக் கண்டுபிடிக்காமல்
நின்றுவிட்டாலோ, ஒருபோதும் நிற்காவிட்டாலோ,
$\Gamma \Sequent \Delta$ தொடக்கத்திலிருந்தே
ஏற்புடையதல்ல.

படிமுறையின் இன்றியமையாத பண்பு ஒன்று: தொடக்கத்
தொடரணி~$\Gamma \Sequent \Delta$ இல் அல்லது
தேடலின்போது உருவான தொடரணியில் உள்ள ஒவ்வொரு
!!{formula} உம் தேவையான அளவுக்குப் ``குறைக்கப்படுவதை''
அது உறுதிசெய்ய வேண்டும். அதாவது, அனுமான விதியைப்
பின்திசையில் பயன்படுத்தும்போது அந்த வாய்ப்பாடு
முதன்மை வாய்ப்பாடாகத் தேர்வாக வேண்டும்; தேவையான
முறைதோறும் இப்படித் தேர்வாக வேண்டும். இப்பண்பை
\emph{நியாயமான தேர்வு} என்போம்.

% SEGMENT OLP-0683-S02
படிமுறை கட்டங்களாகச் செயல்படுகிறது. ஒவ்வொரு
கட்டத்தின் வெளியீடும் தொடரணிகளின் ஒரு மரம்;
அடுத்த கட்டத்தின் வெளியீடு, ஏற்கெனவே அடிகோளாக
இல்லாத ஒவ்வொரு மிக மேலுள்ள தொடரணியிலும்
!!a{formula} ஐக் குறைப்பதால் கிடைக்கிறது.

$0$ ஆம் கட்டத்தில் $\Gamma \Sequent \Delta$
உடன் தொடங்குகிறோம். நியாயமான தேர்வை உறுதிசெய்ய,
$\Gamma$ மற்றும் $\Delta$ விலுள்ள !!{formula}s
ஒவ்வொன்றுக்கும் வெவ்வேறு எண்களைக் குறியீடுகளாக
வழங்குகிறோம். வேறுபட்ட !!{formula}s ஒரே எண்ணைப்
பெறுவதில்லை. எடுத்துக்காட்டாக, இடமிருந்து வலமாக
!!{formula}s ஐ எண்ணலாம். எண்களை மேலெழுத்துகளாக
எழுதுகிறோம். $!A, !B \Sequent !C$ க்கான
!!a{proof} ஐத் தேடுகிறோம் எனில்,
$0$ ஆம் கட்டத்தின் வெளியீடு
$!A^1, !B^2 \Sequent !C^3$ ஆகும். அத்தகைய
குறியிடப்பட்ட தொடரணியில் மேலெழுத்தாக வரும்
மிகப் பெரிய~$n$ அதன் \emph{மிகப் பெரிய குறியீட்டு எண்};
எடுத்துக்காட்டில் அது~$3$.

% SEGMENT OLP-0683-S03
தொடரணியில் அளவையடைகள் இருந்தால்,
$\lexists[x][!A(x)]$ ஐ வலப்புறத்திலோ
$\lforall[x][!A(x)]$ ஐ இடப்புறத்திலோ
குறைக்கும்போது எந்தச் சொற்களைப் பயன்படுத்துவது
என்பதையும் அறிய வேண்டும். $\Gamma$ மற்றும்
$\Delta$ வில் வரும் சார்புக் குறிகளைக் கொண்டு
!!{constant}s ஆன $\Obj c_0$, $\Obj c_1$,
$\Obj c_2$, \dots, இவற்றிலிருந்து கட்டப்படும்
சொற்கள் மட்டுமே தேவை. இச்சொற்கள் அனைத்தும்
ஒரு நிலையான வரிசையில் எண்ணிடப்பட்டுள்ளன என்று
கொள்கிறோம். அப்போதுதான் எடுத்துக்காட்டாக,
``$\Gamma \Sequent \Delta$ இல் இதுவரை
வராத முதல் !!{constant}'' என்று சொல்ல முடியும்.
படிமுறை உருவாக்கும் மரத்தின் ஒவ்வொரு குறியிடப்பட்ட
தொடரணிக்கும் !!{constant}s இன் ஒரு கணம்
இணைக்கப்பட்டுள்ளது. $0$ ஆம் கட்டத் தொடரணிக்கு
இணைக்கப்படும் !!{constant}s இன் கணம், அதில் வரும் எல்லா
!!{constant}s இன் கணம்; எந்த !!{constant}s உம்
இல்லையெனில் அது $\{\Obj c_0\}$.

% SEGMENT OLP-0683-S04
படிமுறை $n$ ஆம் கட்டத்தில், இணைக்கப்பட்ட
!!{constant}s இன் கணங்களோடு குறியிடப்பட்ட
தொடரணிகளின் மரத்தை உருவாக்கியிருக்கட்டும்.
$n+1$ ஆம் கட்டத்தில் மிக மேலுள்ள ஒவ்வொரு
தொடரணிக்கும் பின்வருவதைச் செய்கிறோம்.

ஒரு !!a{formula}~$!A$ தொடரணியின் இருபுறங்களிலும்
வந்தாலோ, இடப்புறத்தில்~$\lfalse$ இருந்தாலோ,
ஒன்றும் செய்யமாட்டோம். அத்தகைய தொடரணிகளுக்கு
\Log{G3c} இல் !!{proof}s உண்டு; அந்த மிக மேலுள்ள
தொடரணியில் மேற்கொண்டு செய்ய வேண்டியது இல்லை.

தொடரணியில் அணுவல்லாத !!{formula} ஏதும்
இல்லையெனில் படிமுறையை நிறுத்துகிறோம்.
$\Gamma \Sequent \Delta$ க்கு நிறுவல் இல்லை.

இல்லாவிட்டால், அணுவல்லாத !!{formula}~$!A$
களில் மிகச் சிறிய குறியீட்டு எண்~$i$ கொண்டதைத்
தேர்ந்தெடுக்கிறோம். அப்போது மிக மேலுள்ள தொடரணி
$!A^i, \Pi \Sequent \Lambda$ அல்லது
$\Pi \Sequent \Lambda, !A^i$ என்ற வடிவில்
இருக்கும். தொடரணியில் ஏற்கெனவே உள்ள மிகப் பெரிய
குறியீட்டு எண்ணைவிட ஒன்றால் பெரிய எண்ணாக~$k$
ஐயும், அதனுடன் இணைக்கப்பட்ட !!{constant}s இன்
கணமாக~$C$ ஐயும் கொள்கிறோம். ஆகவே புதிய
வாய்ப்பாடுகளின் குறியீடுகள் முந்தையவற்றிலிருந்து
வேறுபட்டு அவற்றைவிடப் பெரியவை.

$!A^i$ ஐ ``குறைக்கிறோம்'': $!A$ இன்
!!{main operator} எது என்பதற்கும், $!A^i$
இடப்புறமா வலப்புறமா வருகிறது என்பதற்கும்
ஏற்ற அனுமான விதியின்படி புதிய மிக மேலுள்ள
தொடரணிகளை உருவாக்குகிறோம்.

% SEGMENT OLP-0683-S05
\begin{enumerate}
  \item $!A \ident !B \land !C$ இடப்புறத்தில்
  வருகிறது என்க. $(!B \land !C)^i, \Pi \Sequent \Lambda$
  க்கு மேலாக ஒரு புதிய மிக மேலுள்ள தொடரணியைச்
  சேர்க்கிறோம்:
  \[
  \Axiom$!B^k, !C^{k+1}, \Pi \fCenter \Lambda$
  \RightLabel{\LeftR{\land}}
  \UnaryInf$(!B \land !C)^i, \Pi \fCenter \Lambda$
  \DisplayProof
  \]
  புதிய தொடரணியுடன் இணைக்கப்படும் !!{constant}s
  இன் கணம்~$C$.
  \item $!A \ident !B \land !C$ வலப்புறத்தில்
  வருகிறது என்க. $\Pi \Sequent \Lambda, (!B \land !C)^i$
  க்கு மேலாக இரண்டு புதிய மிக மேலுள்ள
  தொடரணிகளைச் சேர்க்கிறோம்:
  \[
  \Axiom$\Pi \fCenter \Lambda, !B^k$
  \Axiom$\Pi \fCenter \Lambda, !C^k$
  \RightLabel{\RightR{\land}}
  \BinaryInf$\Pi \fCenter \Lambda, (!B \land !C)^i$
  \DisplayProof
  \]
  புதிய தொடரணிகள் ஒவ்வொன்றுடனும் இணைக்கப்படும்
  !!{constant}s இன் கணம்~$C$.

% SEGMENT OLP-0683-S06
  \item $!A \ident \lnot !B$,
  $!A \ident !B \lor !C$, மற்றும்
  $!A \ident !B \lif !C$ ஆகிய நேர்வுகளும்
  இதைப் போன்றவை; பயிற்சிகளாக விடப்படுகின்றன.

  \item $!A \ident \lexists[x][!B(x)]$
  இடப்புறத்தில் வருகிறது என்க.
  $\lexists[x][!B(x)]^i, \Pi \Sequent \Lambda$
  க்கு மேலாக ஒரு புதிய மிக மேலுள்ள
  தொடரணியைச் சேர்க்கிறோம்:
  \[
  \Axiom$!B(c)^k, \Pi \fCenter \Lambda$
  \RightLabel{\LeftR{\lexists}}
  \UnaryInf$\lexists[x][!B(x)]^i, \Pi \fCenter \Lambda$
  \DisplayProof
  \]
  இங்கே~$c$ என்பது $C$ இல் இல்லாத முதல்
  !!{constant}. புதிய தொடரணியுடன் இணைக்கப்படும்
  !!{constant}s இன் கணம் $C \cup \{c\}$.

% SEGMENT OLP-0683-S07
  \item $!A \ident \lexists[x][!B(x)]$
  வலப்புறத்தில் வருகிறது என்க.
  $\Pi \Sequent \Lambda, \lexists[x][!B(x)]$
  க்கு மேலாக ஒரு புதிய மிக மேலுள்ள
  தொடரணியைச் சேர்க்கிறோம்:
  \[
  \Axiom$\Pi \fCenter \Lambda, \lexists[x][!B(x)]^k, !B(t)^{k+1}$
  \RightLabel{\RightR{\lexists}}
  \UnaryInf$\Pi \fCenter \Lambda, \lexists[x][!B(x)]^i$
  \DisplayProof
  \]
  இங்கே~$t$ என்பது $C$ இல் உள்ள !!{constant}s
  மட்டுமே கொண்டு கட்டப்பட்ட சொற்களில்,
  இந்தத் தொடரணிக்குக் கீழே வலப்புறத்தில்
  $\lexists[x][!B(x)]$ ஐக் குறைக்க இதுவரை
  பயன்படுத்தப்படாத முதல் சொல். அத்தகைய சொற்கள்
  அனைத்தும் பயன்படுத்தப்பட்டிருந்தால்
  $\Obj c_0$ ஐ எடுத்துக்கொள்கிறோம்.
  புதிய தொடரணியுடன் இணைக்கப்படும்
  !!{constant}s இன் கணம் $C$.
  \item $!A \ident \lforall[x][!B(x)]$
  என்ற நேர்வுகளும் இதைப் போன்றவை;
  பயிற்சிகளாக விடப்படுகின்றன.
\end{enumerate}

% SEGMENT OLP-0683-S08
$n+1$ ஆம் கட்டத்தில் மிக மேலுள்ள எந்தத்
தொடரணிக்கும் புதிய தொடரணியைச் சேர்க்கவில்லை
எனில், படிமுறை வெற்றிகரமாக நிற்கிறது. அப்போது
$n+1$ ஆம் கட்ட மரத்திலிருந்து குறியீட்டு எண்கள்
அனைத்தையும் நீக்கி, $\Gamma \Sequent \Delta$
க்கான !!a{proof} ஐப் பெறுகிறோம். மிக மேலுள்ள
தொடரணிகள் எல்லாம் $!A, \Pi \Sequent \Lambda, !A$
அல்லது $\lfalse, \Pi \Sequent \Lambda$ என்ற
வடிவிலுள்ளவை. எல்லா அனுமானங்களும் \Log{G3c}
இன் சரியான அனுமானங்கள். கூற்றுக்கணிப்பு விதிகளுக்கு
இது வெளிப்படையானது. ஒவ்வொரு படியிலும் தொடரணியுடன்
இணைக்கப்பட்ட !!{constant}s இன் கணம்~$C$,
அந்தத் தொடரணியில் வரும் எல்லா !!{constant}s
ஐயும் கொண்டுள்ளது என்பதைக் கவனிக்கவும்.
ஆகவே \LeftR{\lexists} மற்றும்
\RightR{\lforall} மூலம் குறைக்கும்போது,
தனிமாறாக் குறி நிபந்தனை நிறைவேறும்:
தனிமாறாக் குறியான~$c$, முடிவுத் தொடரணியுடன்
இணைக்கப்பட்ட~$C$ கணத்தில் இல்லாத
!!a{constant} என்பதால், முடிவிலும்~$c$ வராது.
இதனால் பின்வருவது கிடைக்கிறது:

\begin{prop}
\Log{G3c} க்கான நிறுவல் தேடல் படிமுறை சரியானது:
அது வெற்றிகரமாக நின்றால், தொடக்கத் தொடரணியான
$\Gamma \Sequent \Delta$ க்கு
!!a{proof} உண்டு.
\end{prop}

\end{document}
